\chapter*{\centerline{ТЕРМИНЫ И ОПРЕДЕЛЕНИЯ}}             % Заголовок
\addcontentsline{toc}{chapter}{ТЕРМИНЫ И ОПРЕДЕЛЕНИЯ}  % Добавляем его в оглавление

Большие языковые модели (LLM) -- современные компьютерные модели, основанные на глубоком обучении, которые обрабатывают и генерируют естественный язык, обладая способностью к пониманию и созданию текста на уровне, приближенном к человеческому.

Интеграция эмпатии -- процесс включения эмпатических аспектов и понимания чувств и переживаний других людей в деятельность, взаимодействие или принятие решений.

Пайплайн -- это последовательность этапов, необходимых для тестирования моделей и промптов, начиная с определения базовых параметров и заканчивая анализом и сохранением результатов.

Поведение агентов LLM -- способы взаимодействий и действий, которые могут быть созданы или смоделированы с использованием больших языковых моделей, таких как создание текста, ответ на запросы или генерация содержания в различных контекстах и ситуациях.

Принятие решений -- процесс выбора оптимального варианта действия из множества альтернативных вариантов в условиях неопределенности или ограниченности ресурсов.

Промпт -- в контексте языковых моделей представляет собой инструкцию или описание, предоставляемое пользователем, которое служит руководством для генеративной языковой модели при создании контента.

Теория игр -- математическая дисциплина, изучающая стратегии принятия решений в условиях конфликта и кооперации между рациональными игроками.

Эмпатия -- способность постижения и понимания чувств, мыслей и переживаний других людей, а также способность к сопереживанию и сочувствию.
